\documentclass[../../../include/open-logic-section]{subfiles}

\begin{document}

\olfileid{sth}{z}{infinity-again}	
\olsection{అనంతత్వం}

ఇప్పటికే మన వద్ద ఉన్న
స్వీకృతాలే అనంతంగా అనేక
సమితులు ఉంటాయని
నిర్ధారించడానికి సరిపోతాయి
(ఏదైనా సమితి ఉంటే).
ఎందుకంటే ఒక సమితి
ఉందనుకుందాం; అప్పుడు
\olref[sth][z][sep]{prop:emptyexists}
ప్రకారం $\emptyset$ ఉంటుంది.
ఇప్పుడు ఏ సమితి~$x$కైనా,
\olref[sth][z][pairs]{prop:pairsconsequences}
ప్రకారం $x \cup \{x\}$
ఉంటుంది. దీనిని కొన్ని
సార్లు ప్రయోగిస్తే ఈ
కింది సమితులు లభిస్తాయి:
\begin{enumerate}
	\item[0.] $\emptyset$ %& $0$ & దశ $0$\\
	\item[1.] $ \{\emptyset\}$ %& $1$ & దశ $1$\\
	\item[2.] $\{\emptyset, \{\emptyset\}\}$ %& $2$ & దశ $2$\\
	\item[3.] $\{\emptyset, \{\emptyset\}, \{\emptyset, \{\emptyset\}\}\}$ %&$3$ & దశ $3$\\
	\item[4.] $\{\emptyset, \{\emptyset\}, \{\emptyset, \{\emptyset\}\}, \{\emptyset, \{\emptyset\}, \{\emptyset, \{\emptyset\}\}\}\}$% & $4$ & దశ $4$ 
\end{enumerate}
ఈ సమితులన్నీ ఒకదానికొకటి
భిన్నమని తనిఖీ చేయవచ్చు.

కొన్ని కారణాల వల్ల
లెక్కింపును $0$ నుంచి
మొదలుపెట్టాం. వాటిలో ఒకటి:
``$n$'' అని గుర్తించిన
సమితికి ఖచ్చితంగా $n$
మూలకాలు ఉంటాయని,
(సహజంగా చూస్తే) అది
$n$వ దశలో ఏర్పడుతుందని
తనిఖీ చేయడం కష్టం కాదు.

అయితే దీని వల్ల
\emph{అనంతంగా అనేక} సమితులు
వస్తాయి గానీ, \emph{అనంత
సమితి}---అనంతంగా అనేక
మూలకాలున్న ఒక సమితి---ఉంటుందని
నిర్ధారణ కాదు. ఇది
నిజంగా ముఖ్యం: ఒక
(డెడెకిండ్) అనంత సమితి
లభించకపోతే, డెడెకిండ్
బీజగణితాన్ని నిర్మించలేం.
దానిని సహజ సంఖ్యల
సమితిగా పరిగణించడానికి
మనకు డెడెకిండ్ బీజగణితం
కావాలి. (\olref[sfr][infinite][dedekindsproof]{sec}తో
పోల్చండి.)

ముఖ్యంగా, ఇప్పటివరకు
చెప్పిన స్వీకృతాలు ఏ
అనంత సమితి ఉనికినీ
\emph{నిర్ధారించవు}.
అందువల్ల మరో స్వీకృతాన్ని
నిర్దేశించాలి:

\begin{axiom}[అనంతత్వం]
$\emptyset \in I$ ఉండి,
$x \in I$ అయినప్పుడల్లా
$x \cup \{x\} \in I$
అయ్యే ఒక సమితి $I$
ఉంటుంది.
\begin{align*}
	\exists I( & (\exists o \in I)\forall x\ x \notin o \land {}\\
	& (\forall x \in I)(\exists s \in I)\forall z(z \in s \liff (z \in x \lor z = x)))
\end{align*}
\end{axiom}

అనంతత్వ స్వీకృతం ఇచ్చే
$I$ సమితి డెడెకిండ్
అనంతమని సులభంగా
చూడవచ్చు. దాని
ప్రత్యేక మూలకం
$\emptyset$; $I$పై
ఒకటి-ఒకటి ప్రమేయం
$s(x) = x\cup
\{x\}$తో ఇవ్వబడుతుంది.
ఇప్పుడు \olref[sfr][infinite][dedekind]{thm:DedekindInfiniteAlgebra}
డెడెకిండ్ అనంత సమితి
నుంచి డెడెకిండ్ బీజగణితాన్ని
ఎలా పొందాలో చూపించింది;
దానినే మన సహజ సంఖ్యల
సమితిగా పరిగణిస్తాం.
మరింత కచ్చితంగా:
\begin{defn}\ollabel{defnomega}
అనంతత్వ స్వీకృతం ఇచ్చే
ఏ సమితి $I$నైనా
తీసుకుందాం. $s$
ప్రమేయాన్ని $s(x) = x \cup \{x\}$
అని నిర్వచిద్దాం.
$\omega =
\closureofunder{s}{\emptyset}$
అనుకుందాం. $\omega$లోని
మూలకాలను \emph{సహజ
సంఖ్యలు} అంటాం; $n$
అనేది $\emptyset$కు
$s$ను $n$ సార్లు
వర్తింపజేసిన ఫలితమని
చెబుతాం.
\end{defn}

కొన్ని పేరాల క్రితం
``$n$'' అని గుర్తించిన
సమితిని $n$ సంఖ్య\emph{గా}
పరిగణిస్తామని ఇప్పుడు
వెనక్కి వెళ్లి తనిఖీ
చేయవచ్చు.

ఈ నిర్దేశం ప్రాముఖ్యతను
\olref[sth][z][nat]{sec}లో
చర్చిస్తాం. ప్రస్తుతానికి
దీనితో ఒక సహజమైన
ఫలితాన్ని నిరూపించవచ్చు:

\begin{prop}\ollabel{naturalnumbersarentinfinite}
ఏ సహజ సంఖ్యా డెడెకిండ్
అనంతం కాదు.
\end{prop}

\begin{proof}
నిరూపణ ఆగమనం ద్వారా,
అంటే \olref[sfr][infinite][induction]{thm:dedinfiniteinduction}
ద్వారా. స్పష్టంగా
$0
= \emptyset$ డెడెకిండ్
అనంతం కాదు. ఆగమన
దశలో విపర్యయాన్ని
నిరూపిస్తాం: (అసంభవమైనప్పటికీ)
$s(n)$ డెడెకిండ్
అనంతం అయితే, $n$
కూడా డెడెకిండ్ అనంతం.

$s(n)$ డెడెకిండ్
అనంతం అనుకుందాం;
అంటే $\ran{f}\subsetneq \dom{f} = s(n) = n \cup
\{n\}$ అయ్యే ఏదో
\tetoken{ఒకటి-ఒకటి ప్రమేయం}{!!{injection}}
$f$ ఉంటుంది. రెండు
సందర్భాలను పరిశీలిద్దాం.

\emph{సందర్భం 1: $n \notin \ran{f}$.}
అప్పుడు $\ran{f} \subseteq n$,
మరియు $f(n) \in n$.
$g = \funrestrictionto{f}{n}$
అనుకుందాం; ఇప్పుడు
$\ran{g} =
\ran{f} \setminus \{f(n)\} \subsetneq n = \dom{g}$.
కాబట్టి $n$ డెడెకిండ్
అనంతం.

\emph{సందర్భం 2: $n \in \ran{f}$.}
$m \in \dom{f} \setminus \ran{f}$
ఒకదాన్ని స్థిరపరచి,
$s(n) = n \cup \{n\}$
నిర్వచన సమితిగల $h$
ప్రమేయాన్ని ఇలా
నిర్వచిద్దాం:
\[
h(x) = 
	\begin{cases}
		f(x) & \text{ఒకవేళ }f(x) \neq n\\
		m & \text{ఒకవేళ }f(x)=n
	\end{cases}
\]
అప్పుడు $h$, $f$
అన్నిచోట్లా ఒకేలా
ఉంటాయి; అయితే
$h(f^{-1}(n)) = m
\neq n = f(f^{-1}(n))$
వద్ద మాత్రం భిన్నం.
$f$ ఒక
\tetoken{ఒకటి-ఒకటి ప్రమేయం}{!!a{injection}}
కాబట్టి, $n \notin
\ran{h}$; ఇంకా
$\ran{h} \subsetneq \dom{h} = s(n)$.
ఇప్పుడు సందర్భం 1
వాదనతో $n$ డెడెకిండ్
అనంతమని తేలుతుంది.
\end{proof}

అయినా అనంతత్వ స్వీకృతాన్ని
\emph{ఎలా సమర్థించాలి}
అనే ప్రశ్న మిగిలే ఉంది.
దానికి చిన్న జవాబు:
మనం చెప్పుకుంటున్న
కథకు మరో సూత్రాన్ని
జోడించాలి. అది ఇది:
\begin{enumerate}
	\item[] \stagesinf. ఒక అనంత దశ ఉంది. అంటే, (a) మొదటి దశ కాని, (b) తనకు ముందు కొన్ని దశలున్న, కానీ (c) వెంటనే ముందు దశ లేని ఒక దశ ఉంది.
\end{enumerate}
ఈ సూత్రం నుంచి అనంతత్వ
స్వీకృతం సరళంగా వస్తుంది.
సహజ సంఖ్య $n$,
$n$వ దశలో ఏర్పడుతుందని
మనకు తెలుసు. కాబట్టి
$\omega$ సమితి తొలి
అనంత దశలో ఏర్పడుతుంది.
$\omega$ స్వయంగా
అనంతత్వ స్వీకృతానికి
సాక్ష్యమిస్తుంది.

అయితే ఇది \stagesinf{}ను
ఎలా సమర్థించాలి అనే
ప్రశ్నకే మనల్ని తిరిగి
తీసుకెళ్తుంది.
\stagessucc\ మాదిరిగానే,
సంచిత-దశలవారీ
స్థాయిక్రమంపై మన కథలో
ఇది స్పష్టంగా లేదు.
అంతకంటే ముఖ్యంగా,
సమితులు దశలవారీగా
ఏర్పడే స్థాయిక్రమ భావన
ఒక్కటే ఆ ప్రక్రియలో
\emph{అనంత} దశ ఉండాలని
నిర్బంధించదు. దశలు
సహజ సంఖ్యల మాదిరి
క్రమంలో ఉన్నాయని
భావించడం పూర్తిగా
సంగతంగానే కనిపిస్తుంది.

అయితే దీనివల్ల ఒక
స్పష్టమైన సమస్య
తలెత్తుతుంది.
\olref[sfr][infinite][dedekindsproof]{sec}లో,
(ఆలోచనల) డెడెకిండ్
అనంత సమితి ఉందని
డెడెకిండ్ ఇచ్చిన
``నిరూపణ''ను చూశాం.
అది మీకు అంతగా
సంతృప్తికరంగా అనిపించి
ఉండకపోవచ్చు. కానీ
సమితుల దశలవారీ భావన
(లేదా ``ఆలోచనా
నియమాలు'') \stagesinf{}ను
``మనపై బలవంతంగా''
రుద్దకపోతే, డెడెకిండ్
అనంత సమితి ఉందనే
వాదనకు అంతర్గత
సమర్థన మనకు ఇంకా
ఉండదు.

ఇక్కడ చెప్పాల్సింది
ఇంకా చాలా ఉంది.
అయితే ఒక స్వీకృతాన్ని
(లేదా సూత్రాన్ని)
\emph{సమర్థించాలంటే}
ఏమి అవసరమవుతుందో
ఆలోచించడం ప్రారంభించే
స్థితికి ఇప్పుడు మీరు
చేరారని ఆశిస్తున్నాం.
ఇకముందు \stagesinf{}ను
అంగీకరించి సాగుతాం.

\end{document}
